\documentclass[11pt,article]{article}

% Packages
\usepackage[utf8]{inputenc}
\usepackage[margin=1in]{geometry}
\usepackage{amsmath,amssymb,amsthm}
\usepackage{listings}
\usepackage{xcolor}
\usepackage{hyperref}
\usepackage{enumitem}
\usepackage{graphicx}

% Code listing settings
\lstset{
    basicstyle=\ttfamily\small,
    breaklines=true,
    frame=single,
    language=Python,
    commentstyle=\color{gray},
    keywordstyle=\color{blue},
    stringstyle=\color{red},
    showstringspaces=false,
    numbers=left,
    numberstyle=\tiny\color{gray},
    backgroundcolor=\color{gray!10}
}

% Title
\title{\textbf{TRM Improvement with RL for ARC-2}}
\author{}
\date{}

\begin{document}

\maketitle

\section{RL-style iteration or energy/proximal view}


\subsection{Two equivalent approaches}

\subsubsection{Unrolled policy iteration}

\begin{itemize}
    \item View $y$ as the policy (the current plan/solution) and $z$ as the value/belief summarizing consequences.
    \item The inner update $z \leftarrow f(z, y, x)$ approximates policy evaluation; the outer $y \leftarrow g(y, z, x)$ approximates policy improvement.
    \item \textbf{Bellman-residual loss on $z$}: supervise $f$ to be a contraction toward targets derived from a short rollout/checker.
    \item \textbf{Stability tool}: bound the Lipschitz constant of $f$ via spectral norm / weight normalization $\rightarrow$ guarantees on convergence as $n$ grows.
\end{itemize}

\subsubsection{Learned proximal operator}

\begin{itemize}
    \item Define a learnable energy $E_\theta(y, z; x)$. Update:
    \[
    z_{t+1} \approx \text{prox}_{\alpha R}\big(z_t - \alpha \nabla_z E_\theta(y_t, z_t; x)\big)
    \]
    (prox implemented as a tiny residual MLP with nonexpansiveness enforced).
    
    \item \textbf{Lyapunov}: train $E$ so $\Delta E \le -c\|z_{t+1}-z_t\|^2$. Add a penalty if energy increases.
    
    \item \textbf{Tuning handles}: step size $\alpha$ and nonexpansive prox give us a principled way to pick safe $n,T$ without blowups.
\end{itemize}

\subsection{Practical constraints}

\begin{itemize}
    \item Enforce nonexpansiveness with spectral norm ($\le 1$) and 1-Lipschitz activations (e.g., GroupSort or clamp-tanh for the prox block), or add a Jacobian penalty on $f$.
    \item Add a small gradient-step gate $z \leftarrow z + \beta \Delta z$ with $\beta\in(0,1]$ learned but clipped; this is a learnable step size.
\end{itemize}

\subsection{Why it helps}

We stop guessing $n,T$ and layer counts; the math gives guardrails. It often explains why tiny nets + many recursions beat deep nets in small-data regimes: recursion reuses parameters under a stable dynamical system.



\end{document}